%%%%%%%%%%%%%%%%%%%%%%%%%%%%%%%%%%%%%%%%%%%%%%%%%%%%%%%%%%%%
%% Section 7: Industry Architecture Convergence
%%%%%%%%%%%%%%%%%%%%%%%%%%%%%%%%%%%%%%%%%%%%%%%%%%%%%%%%%%%%
%% \section moved to main.tex
\label{sec:industry}

A systematic analysis of 28 technical reports from 24 organizations (2024--2025) reveals a striking pattern: LLM architectures have converged on a common ``standard block'' at the component level, while diverging rapidly at the system level. This section examines both the convergence and the diversification, and identifies the design principles that emerge from industrial practice.

\subsection{The Converged Standard Block}
\label{subsec:standard_block}

Table~\ref{tab:industry_convergence} summarizes the five de facto standards that have emerged across nearly all production LLMs.

\begin{table}[t]
\centering
\caption{Industrial convergence: architectural choices across 28 production LLMs (2024--2025).}
\label{tab:industry_convergence}
\small
\begin{tabular}{llcc}
\hline
\textbf{Component} & \textbf{Standard Choice} & \textbf{Adoption} & \textbf{Notable Exception} \\
\hline
Normalization & RMSNorm~\cite{zhang2019root} & 26/28 & PanGu-Sigma (LN), MPT (LN) \\
Norm Placement & Pre-Norm & 27/28 & OLMo~2 (Reordered-Norm) \\
Activation & SwiGLU~\cite{shazeer2020glu} & 26/28 & Gemma~2 (GeGLU), Nemotron (Sq.\ ReLU) \\
Position Encoding & RoPE~\cite{su2024roformer} & 27/28 & MPT (ALiBi) \\
Token Mixer & GQA~\cite{ainslie2023gqa} & 22/28 & DeepSeek (MLA), Step-2 (MFA) \\
\hline
\end{tabular}
\end{table}

This convergence is driven by three forces: engineering pragmatism (RMSNorm is computationally cheap, SwiGLU is reliably performant, GQA is inference-efficient); ecosystem effects (Llama~2/3's wide adoption established its architecture as a ``default template,'' driving latecomers like Falcon~3 to migrate from ALiBi to RoPE); and ablation cost (comprehensive architecture ablations at 100B+ scale are prohibitively expensive, incentivizing reuse of validated choices).

\subsection{The Diverging Frontier}
\label{subsec:diverging_frontier}

While component-level choices have converged, system-level design---particularly around Token Mixer strategy, MoE configuration, and position encoding usage---exhibits significant diversification.

\paragraph{Token Mixer diversification.}
Beyond the GQA baseline, four alternative KV-cache optimization strategies have emerged: (1)~\textbf{MLA} (DeepSeek-V2/V3/R1), which jointly compresses KV into a low-dimensional latent vector achieving 93.3\% cache reduction; (2)~\textbf{CLA} (Hunyuan-Large, Yi-Lightning), which shares KV caches across adjacent layers; (3)~\textbf{MFA} (Step-2 mini), which applies matrix factorization to the QK interaction; and (4)~\textbf{heterogeneous attention patterns}---interleaving local/global attention (Gemma~3, 5:1), linear/softmax attention (MiniMax-01, 7:1), or SSM/Transformer layers (Nemotron-3 Nano). QK-Norm, adopted by Gemma~3, Qwen~3, and Llama~4, is becoming a new sub-component standard for attention stability.

\paragraph{MoE becomes mainstream.}
MoE has transitioned from a niche technique (2 of 28 models in 2023) to the dominant paradigm for frontier models ($>$14 of 28 in 2024--2025). Expert granularity spans three regimes: coarse (Mixtral: 8 experts), medium (Qwen~3: 128, ERNIE~4.5: 64), and ultra-fine (DeepSeek-V3: 256 experts with $\sim$2048-dimensional intermediates each). Shared experts remain contentious---DeepSeek-V3 retains one shared expert for cross-domain common patterns, while Qwen~3 removed shared experts entirely. Load balancing has evolved from auxiliary loss~\cite{fedus2022switch} to auxiliary-loss-free routing (DeepSeek-V3) to global-batch balancing (Qwen~3). Activation ratios continue to decline: from 27.8\% (Mixtral) to 5.5\% (DeepSeek-V3) to 4.3\% (Llama~4 Maverick), suggesting that future flagship models may activate only 3--5\% of total parameters per token.

\paragraph{Position encoding innovation.}
RoPE base frequency has grown exponentially---from 10K (original, 2021) to 500K (Llama~3) to 1M (Gemma~3, Qwen~2.5) to 10M (MiniMax-01)---to support longer contexts. More importantly, how RoPE is applied now varies across layers: Decoupled RoPE (DeepSeek) separates positional dimensions from compressed KV, dual-base-frequency RoPE (Gemma~3) assigns different frequencies to local vs.\ global layers, iRoPE (Llama~4) interleaves RoPE and NoPE layers, and half-dimension RoPE (MiniMax-01) applies rotation to only half the head dimensions.

\subsection{The Context Length Arms Race}
\label{subsec:context_race}

\begin{table}[t]
\centering
\caption{Context length evolution and enabling architectural techniques (2023--2025).}
\label{tab:context_length}
\small
\begin{tabular}{lrlll}
\hline
\textbf{Model} & \textbf{Context} & \textbf{Year} & \textbf{Architecture Type} & \textbf{Key Technique} \\
\hline
Llama~2 & 4K & 2023 & Dense & Standard RoPE \\
Mixtral & 32K & 2024 & MoE & SWA + RoPE \\
Llama~3 & 128K & 2024 & Dense & RoPE base=500K \\
DeepSeek-V3 & 128K & 2024 & MoE + MLA & YaRN \\
Hunyuan-Large & 256K & 2024 & MoE + CLA & GQA + CLA \\
Qwen2.5-1M & 1M & 2025 & Dense & Dual Chunk Attention \\
MiniMax-01 & 4M & 2025 & MoE + Linear Attn & Lightning Attention (7:1) \\
Llama~4 Scout & 10M & 2025 & MoE + iRoPE & iRoPE + temp.\ scaling \\
\hline
\end{tabular}
\end{table}

Context length has grown by three orders of magnitude in two years (Table~\ref{tab:context_length}). The enabling techniques form a clear progression: RoPE base-frequency scaling for the 128K frontier, heterogeneous attention patterns (local/global interleaving) for the 1M frontier, and sub-quadratic attention components (linear attention, SSM) combined with position encoding innovations (iRoPE) for the 10M frontier. However, effective utilization of ultra-long contexts---as opposed to mere tolerance---remains an open challenge. Needle-in-a-Haystack evaluations reveal declining retrieval accuracy at extreme lengths, and the quadratic prefill cost of full-attention layers becomes a practical bottleneck regardless of KV cache compression.

\subsection{Outlook: A Three-Layer Design Hierarchy}
\label{subsec:industry_outlook}

The patterns observed across industrial practice suggest a three-layer hierarchy in LLM architecture design:

\begin{enumerate}[leftmargin=*]
  \item \textbf{Converged foundations} (limited innovation space): RMSNorm, SwiGLU, RoPE, Pre-Norm residual. These choices are unlikely to change in the near term; innovations build atop them (e.g., QK-Norm on top of RMSNorm) rather than replacing them.

  \item \textbf{Active diversification} (current innovation frontier): Token Mixer strategy (GQA vs.\ MLA vs.\ CLA vs.\ MFA vs.\ hybrid attention), Channel Mixer design (Dense vs.\ MoE with varying granularity, routing, and shared-expert configurations), and position encoding usage (base frequency, per-layer differentiation, partial application). This layer exhibits the greatest variation across organizations.

  \item \textbf{Emerging directions} (next-generation frontier): fully heterogeneous layer design (each layer or layer group with distinct attention type, FFN width, expert configuration, and PE strategy), dynamic architectures (input-adaptive depth via early exit, dynamic expert count, conditional computation), and training-objective architecturalization (multi-token prediction~\cite{liu2024deepseek} embedded as an architectural component rather than a post-training addition).
\end{enumerate}

The architectural innovation focus is thus shifting from intra-block component selection (largely settled) to inter-block composition and heterogeneity (actively contested). Dense models retain competitiveness at smaller scales ($<$30B) through data quality~\cite{abdin2024phi4} and overtraining, but MoE will dominate frontier models. The central open question is whether architectures will fully converge---with all organizations eventually adopting a single optimal template---or whether the system-level design space is sufficiently rich to sustain persistent diversification driven by different hardware, data, and application constraints.
